\section{Tuning}
\label{sec:tuning}

\subsection{Search space and protocol}

Only parameters the solver actually implements take part in the search:
pricing scheme and its block size or candidate limit, entering rule, forest
update mode with its rebuild interval and fraction, canonicalisation,
degeneracy trigger, vertex ordering of the condensation, transitive reduction and
the initial basis. Conditional parameters are sampled only when their guard
holds, so a candidate never carries a block size it cannot use, and the
constraint that Bland's rule requires full pricing mirrors the solver's own
validation. Numerical tolerances are deliberately excluded: relaxing them would
buy time by solving a different problem.

Candidates are drawn by reproducible random search from a fixed seed, after a
small list of a-priori plausible configurations. Each candidate is evaluated on
a deterministic subsample of the training manifest under a fixed per-run limit,
and ranked lexicographically by number of incorrect results, number of instances
solved and verified within the limits, PAR-2 time, and peak memory. Every
attempt, including rejected ones and the reason for rejection, is appended to a
history file, so an interrupted search resumes without repeating work and the
total cost of tuning --- not only the cost of the winning configuration --- is
recorded.

\generated{table_tuning_forest-v2}

A first search was run before the parameters that turn out to matter existed in
the solver, and its attempts differed by less than three per cent: it is recorded
in \file{experiments/tuning/} but it selected nothing and is not reported here.
The search shown above is the one whose winner was frozen.

\subsection{Frozen configuration}

The configuration selected on training and validation data is written to
\file{configs/frozen/} together with its date, the split it was chosen on, the
search space and the rationale, and only then is the test manifest measured. The
test set is never used to adjust a default.
